\chapter{سپری از اعداد؛ طراحی بستر تبادل داده و رمزگذاری همریخت}
\label{chap:cryptography}
آزمایش فصل قبل نشان داد که می‌توان دادۀ آموزش را میان کلاینت‌ها تقسیم و احتمال‌های خروجی مدل‌های محلی را تجمیع کرد. این جداسازی به‌خودی‌خود امنیت کانال یا محرمانگی خروجی منفرد را تضمین نمی‌کند. برای تحلیل لایۀ تبادل داده، باید میان حفاظت از متن ورودی، مدیریت کلید، اصالت پیام و محاسبه بر متن رمز تمایز گذاشت. در این فصل، ابتدا رمزگذاری احرازاصالت‌شده و تسهیم راز، سپس پنج طرح همریخت و در پایان نتایج بنچمارک موجود بررسی می‌شوند.

در معماری مفهومی، سرور یک متن پرسش مشترک را برای پیش‌بینی در اختیار کلاینت‌ها قرار می‌دهد. هر کلاینت از مدل محلی خود احتمال کلاس مثبت را محاسبه می‌کند. تجمیع جمعی این احتمال‌ها به طرحی با همریختی جمعی نیاز دارد؛ همریختی ضربی به‌تنهایی میانگین وزنی مورد نظر را محاسبه نمی‌کند. همچنین، انتشار کلید عمومی باید به هویت معتبر متصل باشد؛ عمومی بودن کلید، نیاز به اصالت‌سنجی آن را از میان نمی‌برد.

\section{رمزگذاری متن ورودی با \en{AES-GCM}}
\en{AES} یک رمز بلوکی متقارن است؛ حالت \en{GCM} آن رمزگذاری احرازاصالت‌شده با دادۀ مرتبط را فراهم می‌کند \cite{nist-gcm}. کلید نشست ۲۵۶بیتی $K$ با مولد تصادفی امن سیستم‌عامل تولید می‌شود. برای هر پیام، مقدار یک‌بارمصرف $N$، معمولاً با طول ۹۶ بیت، باید تحت همان کلید تکرار نشود. رابطۀ رمزگذاری به‌صورت
\begin{equation}
(C,T)=\Enc_{\mathrm{AES\text{-}GCM}}(K,N,M,A)
\label{eq:aes-gcm}
\end{equation}
است. $M$ متنِ تبدیل‌شده به بایت، $A$ دادۀ مرتبطِ احرازاصالت‌شده، $C$ متن رمز و $T$ برچسب احراز اصالت است. دادۀ مرتبط می‌تواند شناسۀ پروتکل، دور یا نسخه را شامل شود؛ این داده رمز نمی‌شود، اما تغییر آن باید اعتبار برچسب را از بین ببرد.

در رمزگشایی، گیرنده باید پیش از استفاده از متن، برچسب را اعتبارسنجی کند:
\begin{equation}
\Dec_{\mathrm{AES\text{-}GCM}}(K,N,C,A,T)=
\begin{cases}M,&\text{برچسب معتبر باشد},\\\bot,&\text{در غیر این صورت}.
\end{cases}
\end{equation}
اعتبارسنجی ناموفق باید به رد پیام بینجامد. تصادفی بودن \en{nonce} نیاز به کنترل احتمال برخورد و سیاست تعویض کلید را حذف نمی‌کند. استفادۀ مجدد از یک \en{nonce} با همان کلید می‌تواند امنیت \en{GCM} را به‌شدت تضعیف کند. در اینجا «اصالت» به معنای معتبر بودن پیام نسبت به دارندۀ کلید مشترک است و با هویت غیرقابل‌انکار یک فرستندۀ مشخص یکسان نیست.

\subsection{ضرورت و مبانی تسهیم راز شامیر}
طرح شامیر امکان تعریف سیاست آستانه‌ای برای دسترسی به یک راز را فراهم می‌کند \cite{shamir}. در طرح $(t,n)$، راز به $n$ سهم تقسیم می‌شود؛ هر $t$ سهم معتبر برای بازسازی کافی است و کمتر از $t$ سهم، در مدل ایده‌آل، اطلاعاتی دربارۀ راز نمی‌دهد. انتخاب $t=n$ همکاری تمام اعضا را لازم می‌کند، اما تحمل خروج یا خرابی یک عضو را کاهش می‌دهد. انتخاب $t<n$ دسترس‌پذیری بیشتری فراهم می‌کند و هم‌زمان مجموعۀ کوچک‌تری را برای بازسازی مجاز می‌سازد.

راز صحیح $s$ باید در میدان $\F_p$ قرار گیرد. برای بازنمایی یک کلید ۲۵۶بیتی، انتخاب عدد اول $p=2^{521}-1$ ظرفیت کافی دارد. ضرایب $a_1,\ldots,a_{t-1}$ مستقل و یکنواخت در میدان انتخاب می‌شوند و
\begin{equation}
f(X)=s+\sum_{j=1}^{t-1}a_jX^j\pmod p.
\label{eq:shamir-polynomial}
\end{equation}
برای نقاط متمایز و ناصفر $x_i\in\F_p$، سهم عضو $i$ برابر $(x_i,y_i)$ با $y_i=f(x_i)$ است. می‌توان $x_i=i$ را برای $1\le i\le n<p$ برگزید. عمومی بودن مختصۀ $x_i$ مشکلی ایجاد نمی‌کند؛ مؤلفۀ $y_i$ باید محرمانه و اصالت‌سنجی‌شده به گیرنده برسد.

\subsection{بازسازی کلید با درونیابی لاگرانژ}
اگر $I$ مجموعۀ $t$ عضو مشارکت‌کننده باشد، راز از ارزیابی چندجمله‌ای در صفر بازیابی می‌شود:
\begin{align}
s=f(0)&=\sum_{i\in I}y_i\ell_i(0)\pmod p,\\
\ell_i(0)&=\prod_{\substack{j\in I\\j\ne i}}
\frac{-x_j}{x_i-x_j}\pmod p.
\label{eq:lagrange-reconstruction}
\end{align}
تقسیم در این رابطه به معنای ضرب در وارون پیمانه‌ای است. تمایز نقاط تضمین می‌کند مخرج‌ها ناصفر باشند. یک چندجمله‌ای با درجۀ حداکثر $t-1$ با $t$ نقطه به‌طور یکتا تعیین می‌شود؛ بنابراین بازسازی درست، همان جملۀ ثابت اولیه را به دست می‌دهد.

برای نمونه، در پیکربندی $(t,n)=(5,5)$، چهار سهم برای تعیین راز کافی نیست؛ در پیکربندی $(3,5)$، هر سه سهم معتبر امکان بازسازی را دارند. این تفاوت صرفاً یک انتخاب عددی نیست؛ سیاست اعتماد و دسترس‌پذیری سامانه را تعریف می‌کند. افزودن سهم‌های بیشتر از حد آستانه، بدون اعتبارسنجی سهم‌ها، از خرابکاری یک عضو جلوگیری نمی‌کند؛ تسهیم راز قابل‌راستی‌آزمایی مسئله‌ای جداگانه است.

\subsection{ویژگی امنیتی و حدود آن}
در مدل شامیر با ضرایب یکنواخت، برای هر مقدار فرضی راز و هر مجموعۀ کمتر از $t$ سهم، ضرایبی سازگار وجود دارد. ازاین‌رو مشاهدۀ آن سهم‌ها توزیع راز را تغییر نمی‌دهد؛ این یک ویژگی اطلاعاتی است و به محدود بودن توان محاسباتی مهاجم وابسته نیست. در مقابل، امنیت \en{AES-GCM} محاسباتی و وابسته به پارامترها و استفادۀ صحیح از کلید و \en{nonce} است.

تسهیم راز به این معنا نیست که کلید هیچ‌گاه به‌طور کامل در اختیار هیچ موجودیتی قرار نمی‌گیرد. سرورِ تولیدکننده از ابتدا کلید را می‌داند و گیرندۀ مجاز پس از بازسازی نیز کلید کامل را در اختیار دارد. اگر هدف، جلوگیری از افشای کلید حتی هنگام استفاده باشد، به پروتکل محاسبۀ آستانه‌ای یا سخت‌افزار و مدل اعتماد دیگری نیاز است. همچنین، ارسال تمام سهم‌ها روی یک مسیر قابل‌شنود، حفاظت آستانه‌ای را در برابر شنودگر همان مسیر بی‌اثر می‌کند.

\subsection{جریان کامل پروتکل و مدل تهدید}
جریان مفهومی از دریافت رشتۀ ورودی، تبدیل \en{UTF-8}، تولید کلید و \en{nonce}، رمزگذاری و تولید برچسب، ایجاد سهم‌ها، توزیع امن سهم‌ها، گردآوری تعداد کافی، بازسازی کلید، بررسی برچسب و تبدیل بایت‌های معتبر به متن تشکیل می‌شود. در تمام مراحل، قالب پیام، شناسۀ دور و ارتباط میان متن رمز و سهم‌های همان نشست باید روشن باشند.

\begin{figure}[htbp]
\centering\begin{tikzpicture}[
  node distance=5mm,
  stage/.style={draw=teal!60!black,fill=teal!3,rounded corners=3pt,
    inner sep=6pt,font=\small},
  flow/.style={-{Latex[length=2.2mm]},thick,draw=teal!60!black}]
\node[stage] (text) {\diagramtext{متن ورودی و تبدیل برگشت‌پذیر به بایت با \en{UTF-8}}};
\node[stage,below=of text] (aes) {\diagramtext{تولید کلید تصادفی $K$ و \en{nonce} یکتا برای همان کلید\\
  $(C,T)=\operatorname{AES\text{-}GCM.Enc}(K,IV,M,AAD)$}};
\node[stage,below=of aes] (share) {\diagramtext{تسهیم کلید با آستانۀ $t$ در میدان $\F_p$ با $p>2^{256}$\\
  $f(z)=K+\sum_{j=1}^{t-1}a_jz^j,\quad a_j\overset{\$}{\leftarrow}\F_p$\\
  سهم کلاینت $i$: $(i,f(i))$؛ شناسه‌ها متمایز و ناصفرند}};
\node[stage,below=of share] (send) {\diagramtext{توزیع هر سهم از کانال امن به صاحب آن\\
  نگهداری متن رمز، برچسب و پارامترهای نشست؛ حفاظت از ارتباط آن‌ها}};
\node[stage,below=of send] (recover) {\diagramtext{گردآوری دست‌کم $t$ سهم و درون‌یابی لاگرانژ\\
  $K=f(0)=\sum_{i\in I} f(i)\prod_{j\in I\setminus\{i\}}\frac{-j}{i-j}\pmod p$}};
\node[stage,below=of recover] (verify) {\diagramtext{بررسی برچسب \en{GCM} پیش از تحویل هر متن آشکار\\
  شکست بررسی: رد پیام؛ موفقیت: رمزگشایی معتبر و بازگشایی \en{UTF-8}}};
\node[stage,below=of verify] (trust) {\diagramtext{تولیدکنندۀ کلید در ابتدا $K$ را می‌داند\\
  تسهیم راز، آستانۀ بازسازی را کنترل می‌کند؛ تضمین حذف نسخۀ کلید نیست}};
\foreach \a/\b in {text/aes,aes/share,share/send,send/recover,recover/verify,verify/trust}
  \draw[flow] (\a) -- (\b);
\end{tikzpicture}

\caption{نمای مفهومی حفاظت از متن ورودی با رمزگذاری متقارن و تسهیم راز؛ کانال امن سهم‌ها و کنترل آستانه از پیش‌شرط‌های طرح‌اند}
\label{fig:aes-shamir}
\end{figure}

شنودگر شبکه بدون کلید یا تعداد کافی سهم نباید به متن برسد. مهاجم فعال ممکن است پیام، سهم یا شناسۀ نشست را تغییر دهد؛ برچسب \en{GCM} فقط بخش‌های تحت پوشش خود را کنترل می‌کند و برای سهم‌های نامعتبر به کنترل جداگانه نیاز است. تبانی $t$ عضو برای بازسازی، نقض طرح نیست؛ دقیقاً همان سیاست مجاز طرح آستانه‌ای است. افشای کلید در حافظۀ گیرنده، نشت کانال جانبی، بازیابی پس از خرابی و ثبت‌نشدن مناسب \en{nonce} نیز خارج از تضمین جبری شامیر هستند.

ارزیابی سادۀ چندجمله‌ای برای $n$ سهم، $O(nt)$ عملیات میدان و بازسازی مستقیم با $t$ سهم، $O(t^2)$ عملیات میدان دارد. رمزگذاری متقارن نسبت به طول پیام خطی است. این مرتبه‌ها، تحلیل الگوریتمی‌اند؛ زمان دیواری واقعی به کتابخانه، طول پیام، سخت‌افزار و هزینۀ شبکه وابسته است. در نتایج عددی این پایان‌نامه، زمان رمزنگاری همریخت و زمان آموزش مدل به‌طور جداگانه گزارش می‌شوند.

\section{ارسال خروجی کلاینت‌ها با رمزگذاری همریخت}
فرض کنید تابع مورد نظر، مجموع یا میانگین وزنی مقادیر عددی باشد. طرح همریخت جمعی اجازه می‌دهد از متن‌های رمز، متن رمزِ مجموع ساخته شود. طرح همریخت ضربی، متن رمزِ حاصل‌ضرب را تولید می‌کند. این دو خاصیت نباید با هم جابه‌جا شوند. علاوه بر این، انجام عملیات روی متن رمز، به‌خودی‌خود احراز هویت، کنترل وزن، جلوگیری از بازپخش یا صحت ورودی اعلام‌شده را تأمین نمی‌کند.

\subsection{رمزگذاری \en{ElGamal}}
طرح الجمال در سال ۱۹۸۵ بر پایۀ ساختار گروهی و لگاریتم گسسته معرفی شد \cite{ref178}. در گروه چرخه‌ای مناسب $G=\langle g\rangle$ با مرتبۀ $q$، کلید خصوصی $x$ به‌طور تصادفی انتخاب و کلید عمومی $h=g^x$ محاسبه می‌شود. پیام باید عضوی از $G$ باشد. برای تصادف تازۀ $k$ داریم:
\begin{align}
PK&=(G,q,g,h),&SK&=x,\\
\Enc(m;k)&=(c_1,c_2)=(g^k,mh^k),&
\Dec(c_1,c_2)&=c_2(c_1^x)^{-1}.
\label{eq:elgamal}
\end{align}
عملیات در گروه انجام می‌شوند؛ در نمایش پیمانه‌ای، ضرب و وارون در پیمانۀ عدد اول مربوط محاسبه می‌شوند. درستی مستقیم است، زیرا $h^k=g^{xk}$ و $c_1^x=g^{kx}$، پس عامل تصادفی حذف می‌شود.

اگر $C_1=\Enc(m_1;k_1)$ و $C_2=\Enc(m_2;k_2)$ باشند، ضرب مؤلفه‌ای آن‌ها می‌دهد:
\begin{equation}
C_1\odot C_2=(g^{k_1+k_2},m_1m_2h^{k_1+k_2})
=\Enc(m_1m_2;k_1+k_2).
\end{equation}
بنابراین خاصیت اصلی این صورت از الجمال، ضربی است. رمز کردن توانِ یک مولد می‌تواند صورت جمعی برای مقادیر محدود ایجاد کند، اما رمزگشایی آن به بازیابی توان و محدودیت‌های اضافی نیاز دارد؛ این تبدیل، همان طرح ضربی بنچمارک حاضر نیست.

امنیت معنایی صورت استاندارد الجمال به فرض مناسب \en{DDH} در گروه انتخاب‌شده وابسته است. صرف بزرگ بودن عدد اول یا انتخاب یک مولد، این فرض را برای هر گروه دلخواه برقرار نمی‌کند. نسخۀ پایه نیز تغییرپذیر است و حفاظت در برابر متن رمز انتخاب‌شده را خودکار فراهم نمی‌کند. تولید گروه، اعتبارسنجی اعضا، تصادف تازه و حفاظت از کلید خصوصی برای کاربرد عملی ضروری‌اند.
\begin{figure}[htbp]
\centering\begin{tikzpicture}[
  node distance=6mm,
  stage/.style={draw=blue!50!black,fill=blue!3,rounded corners=3pt,
    inner sep=7pt,font=\small},
  flow/.style={-{Latex[length=2.2mm]},thick,draw=blue!50!black}]
\node[stage] (key) {\diagramtext{گروه دوری معتبر $G=\langle g\rangle$ با مرتبۀ $q$\\
  کلید خصوصی $x$؛ کلید عمومی $y=g^x$؛ توان‌ها در $\Z_q$}};
\node[stage,below=of key] (enc) {\diagramtext{پیام گروهی $m\in G$ و توان تصادفی تازۀ $k$\\
  $c=(c_1,c_2)=(g^k,my^k)$}};
\node[stage,below=of enc] (agg) {\diagramtext{خاصیت ضربی با ضرب مؤلفه‌به‌مؤلفه\\
  $c\otimes c'=(c_1c'_1,c_2c'_2)$\\
  این عملیات متناظر با ضرب پیام‌هاست}};
\node[stage,below=of agg] (dec) {\diagramtext{رمزگشایی با وارونۀ گروهی\\
  $m=c_2(c_1^x)^{-1}$\\
  $\Dec(c\otimes c')=mm'$؛ عملیات در گروه $G$ انجام می‌شوند}};
\node[stage,below=of dec] (security) {\diagramtext{امنیت معنایی مستلزم فرض مناسب \en{DDH} در گروه است\\
  طرح پایه تغییرپذیر است و احراز اصالت پیام را فراهم نمی‌کند}};
\foreach \a/\b in {key/enc,enc/agg,agg/dec,dec/security}\draw[flow] (\a) -- (\b);
\end{tikzpicture}

\caption{طرح مفهومی رمزگذاری و رمزگشایی الجمال}
\label{fig:elgamal}
\end{figure}

\subsection{رمزگذاری \en{Paillier}}
پایۀ طرح پاییه، ساختار باقیمانده‌های مرکب است \cite{ref179}. دو عدد اول متمایز $p,q$ انتخاب، $n=pq$ و $\lambda=\lcm(p-1,q-1)$ محاسبه می‌شوند. تابع $L(u)=(u-1)/n$ برای $u\equiv1\pmod n$ تعریف می‌شود. عنصر $g\in\Z_{n^2}^{*}$ باید به‌گونه‌ای انتخاب شود که مقدار زیر وجود داشته باشد:
\begin{equation}
\mu=\bigl[L(g^{\lambda}\bmod n^2)\bigr]^{-1}\pmod n.
\end{equation}
کلید عمومی $(n,g)$ و کلید خصوصی $(\lambda,\mu)$ است. برای $m\in\Z_n$ و تصادف مستقل $u\in\Z_n^*$، رابطه‌های رمزگذاری و رمزگشایی عبارت‌اند از
\begin{align}
c&=g^m u^n\pmod{n^2},\\
m&=L(c^{\lambda}\bmod n^2)\mu\pmod n.
\label{eq:paillier}
\end{align}
انتخاب متداول $g=n+1$، در صورت برقرار بودن شرایط وارون‌پذیری، محاسبات را ساده می‌کند. نقش تصادف $u$ این است که رمزگذاری یک پیام ثابت هر بار همان متن رمز را تولید نکند.

درستی از حذف عامل تصادفی پس از توان‌رسانی و خطی شدن توان $g$ در تابع $L$ حاصل می‌شود. اگر $c_i=g^{m_i}u_i^n$ باشد، داریم:
\begin{align}
c_1c_2&=g^{m_1+m_2}(u_1u_2)^n\pmod{n^2},\\
\Dec(c_1c_2)&=m_1+m_2\pmod n,&
\Dec(c^w)&=wm\pmod n.
\label{eq:paillier-homomorphism}
\end{align}
توان $w$ در دامنۀ رمز معادل ضرب در عدد آشکار $w$ در دامنۀ پیام است؛ به این ترتیب مجموع وزن‌دار بدون رمزگشایی تک‌تک پیام‌ها ساخته می‌شود. اگر هدف، مجموع صحیحِ معمول باشد، نه مجموع پیمانه‌ای، باید پیش از عملیات از نبود سرریز اطمینان حاصل شود.

پاییه برای تجمیع جمعی مناسب است، اما امضای فرستنده، سیاست وزن و کنترل بازپخش را فراهم نمی‌کند. نسخۀ پایه با وجود امنیت تحت فرض مربوط، در برابر هر نوع مهاجم فعال مصون نیست. تبدیل‌های عمومیِ مقاوم در برابر متن رمز انتخاب‌شده نیز لزوماً خاصیت همریختی مورد نیاز را حفظ نمی‌کنند؛ افزودن یک لایۀ امنیتی باید با هدف جبری طرح سازگار بررسی شود.
\begin{figure}[htbp]
\centering\begin{tikzpicture}[
  node distance=6mm,
  stage/.style={draw=blue!50!black,fill=blue!3,rounded corners=3pt,
    inner sep=7pt,font=\small},
  flow/.style={-{Latex[length=2.2mm]},thick,draw=blue!50!black}]
\node[stage] (key) {\diagramtext{تولید کلید پاییه با پارامترهای معتبر\\
  $n=pq,\quad\lambda=\lcm(p-1,q-1)$\\
  $\mu=[L(g^\lambda\bmod n^2)]^{-1}\bmod n$}};
\node[stage,below=of key] (enc) {\diagramtext{رمزگذاری پیام با تصادف تازه\\
  $m\in\Z_n,\quad u\in\Z_n^*,\quad c=g^m u^n\bmod n^2$}};
\node[stage,below=of enc] (agg) {\diagramtext{تجمیع جمعی و وزن‌دار در دامنۀ رمز\\
  $C=\prod_i c_i^{w_i}\bmod n^2$}};
\node[stage,below=of agg] (dec) {\diagramtext{رمزگشایی با کلید خصوصی\\
  $L(v)=(v-1)/n,\quad M=L(C^\lambda\bmod n^2)\mu\bmod n$\\
  $M=\sum_iw_im_i\pmod n$؛ برای مجموع صحیح، نبود سرریز لازم است}};
\foreach \a/\b in {key/enc,enc/agg,agg/dec}\draw[flow] (\a) -- (\b);
\end{tikzpicture}

\caption{نمای مفهومی پاییه؛ محاسبات متن رمز در پیمانۀ $n^2$ انجام می‌شوند}
\label{fig:paillier}
\end{figure}

\subsection{رمزگذاری \en{Benaloh}}
طرح بنالوح فضای پیام بلوکی $\Z_r$ را به کار می‌گیرد \cite{ref180}. برای صورت متداول، اعداد اول به‌گونه‌ای انتخاب می‌شوند که
\begin{equation}
r\mid(p-1),\qquad \gcd\left(r,\frac{p-1}{r}\right)=1,
\qquad \gcd(r,q-1)=1.
\end{equation}
تعویض نام $p$ و $q$ این ساختار را تغییر نمی‌دهد؛ پیاده‌سازی مخزن، شرط بخش‌پذیری را بر عامل دیگر اعمال می‌کند. با $n=pq$ و $\varphi(n)=(p-1)(q-1)$، باید $y\in\Z_n^*$ به‌گونه‌ای انتخاب شود که $a=y^{\varphi(n)/r}\bmod n$ دارای مرتبۀ دقیق $r$ باشد. وقتی $r$ اول است، شرط $a\ne1$ همراه با شرایط پارامترها کافی است. برای $r$ مرکب، بررسی عامل‌های اولِ مرتبه نیز لازم است.

کلید عمومی $(n,y,r)$ و مادۀ خصوصی شامل عوامل $n$ است. رمزگذاری پیام $m\in\{0,\ldots,r-1\}$ با تصادف $u\in\Z_n^*$ چنین تعریف می‌شود:
\begin{equation}
c=y^m u^r\pmod n.
\label{eq:benaloh}
\end{equation}
برای رمزگشایی، مقدار $b=c^{\varphi(n)/r}\bmod n$ محاسبه می‌شود. چون $u^{\varphi(n)}\equiv1\pmod n$، داریم
\begin{equation}
b=a^m\pmod n.
\label{eq:benaloh-log}
\end{equation}
بازیابی $m$ یک لگاریتم گسسته در زیرگروه با مرتبۀ معلوم $r$ است. برای فضای پیام کوچک، جدول کامل قابل استفاده است؛ پیاده‌سازی حاضر از \en{baby-step/giant-step} با هزینۀ زمانی و حافظۀ مرتبۀ $O(\sqrt r)$ استفاده می‌کند. کوچک بودن فضای پیام یک انتخاب محاسباتی مهم است و افزایش آن بدون اندازه‌گیری هزینۀ رمزگشایی مجاز نیست.

خاصیت جمعی از ضرب متن‌های رمز نتیجه می‌شود:
\begin{align}
c_1c_2&=y^{m_1+m_2}(u_1u_2)^r\pmod n,\\
\Dec(c_1c_2)&=m_1+m_2\pmod r,&
\Dec(c^w)&=wm\pmod r.
\label{eq:benaloh-homomorphism}
\end{align}
در بنچمارک موجود، $r=10000019$ است. این مقدار برای نمایش چهاررقم اعشار و مجموع وزن محدود مناسب است، اما برای وزن‌های صحیحِ برابر تعداد تمام نمونه‌های محلی فصل قبل کافی نیست. شرط ظرفیت در رابطۀ~\eqref{eq:fl-benaloh-capacity} و فصل بعد صریحاً بررسی می‌شود.
\begin{figure}[htbp]
\centering\begin{tikzpicture}[
  node distance=6mm,
  stage/.style={draw=teal!60!black,fill=teal!3,rounded corners=3pt,
    inner sep=7pt,font=\small},
  flow/.style={-{Latex[length=2.2mm]},thick,draw=teal!60!black}]
\node[stage] (key) {\diagramtext{تولید کلید بنالوح با شرایط بخش‌پذیری و هم‌اولی متن\\
  $n=pq,\quad a=y^{\varphi(n)/r}\bmod n,\quad\operatorname{ord}(a)=r$\\
  کلید عمومی $(n,y,r)$؛ فضای پیام $\Z_r$}};
\node[stage,below=of key] (enc) {\diagramtext{رمزگذاری با تصادف تازه\\
  $0\le m<r,\quad u\in\Z_n^*,\quad c=y^m u^r\bmod n$}};
\node[stage,below=of enc] (agg) {\diagramtext{تجمیع وزن‌دار در پیمانۀ $n$\\
  $C=\prod_i c_i^{w_i}\bmod n$}};
\node[stage,below=of agg] (dec) {\diagramtext{رمزگشایی در زیرگروه مرتبۀ $r$\\
  $b=C^{\varphi(n)/r}\bmod n,\quad a^M=b\pmod n$\\
  بازیابی $M=\sum_iw_im_i\pmod r$ با لگاریتم گسستۀ محدود}};
\node[stage,below=of dec] (limit) {\diagramtext{شرط بازیابی مجموع صحیح: $\sum_iw_im_i<r$\\
  در نمایش $0\le m_i\le S$، شرط کافی $S\sum_iw_i<r$ است}};
\foreach \a/\b in {key/enc,enc/agg,agg/dec,dec/limit}\draw[flow] (\a) -- (\b);
\end{tikzpicture}

\caption{نمای مفهومی بنالوح؛ مرتبۀ فضای پیام و روش لگاریتم گسسته بر هزینۀ رمزگشایی اثر دارند}
\label{fig:benaloh}
\end{figure}

\subsection{تعمیم \en{Damgård\lr{--}Jurik} و تخصص \texorpdfstring{$s=1$}{s=1}}
دامگارد\lr{--}یوریک فضای پیام پاییه را با پارامتر $s\ge1$ گسترش می‌دهد \cite{ref181}. در صورت متداول، پیام در $\Z_{n^s}$ و متن رمز در $\Z_{n^{s+1}}^*$ قرار دارد:
\begin{equation}
c=(1+n)^m u^{n^s}\pmod{n^{s+1}}.
\end{equation}
خاصیت جمعی در فضای بزرگ‌تر حفظ می‌شود: $\Dec(c_1c_2)=m_1+m_2\pmod{n^s}$. با این حال، تعمیم مستقیمِ رمزگشایی پاییه با همان $L(u)=(u-1)/n$ برای هر $s$ درست نیست. جمله‌های مرتبۀ بالاترِ بسط دوجمله‌ای باید لحاظ شوند:
\begin{equation}
(1+n)^a\equiv\sum_{j=0}^{s}\binom{a}{j}n^j\pmod{n^{s+1}}.
\end{equation}
اگر $L_s$ وارون نگاشت $a\mapsto(1+n)^a$ در زیرگروه مناسب و پیمانۀ $n^s$ باشد، رمزگشایی را در شرایط معتبر پارامترها می‌توان به‌صورت $m=L_s(c^{\lambda})\lambda^{-1}\pmod{n^s}$ نوشت. محاسبۀ $L_s$ به استخراج مرحله‌ای رقم‌های توان نیاز دارد؛ این نماد همان تقسیم سادۀ تابع $L$ در پاییه نیست.

برای $s=1$، جمله‌های بالاتر حذف و رمزگشایی به صورت پاییه کاهش می‌یابد. کد و بنچمارک فعلی مخزن صریحاً همین تخصص $s=1$ را پشتیبانی می‌کنند. بنابراین زمان‌های گزارش‌شده، شاهد اجرای نسخۀ عمومی برای $s>1$ نیستند. مزیت نظری فضای پیامِ قابل‌تنظیم، بدون پیاده‌سازی و آزمون نسخۀ عمومی نباید به‌عنوان قابلیت اندازه‌گیری‌شده معرفی شود.

\subsection{رمزگذاری \en{RSA}}
طرح \en{RSA} با انتخاب $n=pq$ و توان‌های $e,d$ که $ed\equiv1\pmod{\varphi(n)}$ باشند تعریف می‌شود \cite{ref182}. توان $e$ باید نسبت به $\varphi(n)$ اول باشد. در صورت پایه و بدون لایۀ رمزگذاری احتمالی داریم:
\begin{equation}
PK=(n,e),\quad SK=(n,d),\qquad c=m^e\pmod n,\qquad m=c^d\pmod n.
\label{eq:rsa}
\end{equation}
برای پیام‌های نسبتاً اول با $n$، قضیۀ اویلر درستی را نشان می‌دهد؛ برای حالت کلی نیز استدلال پیمانۀ $p$ و $q$ و قضیۀ باقیماندۀ چینی نتیجه را برقرار می‌کند. این تفکیک مانع کاربرد نادرست قضیۀ اویلر بر هر پیام دلخواه می‌شود.

نسخۀ پایه همریختی ضربی دارد:
\begin{equation}
\Dec(c_1c_2)=m_1m_2\pmod n,\qquad
\Dec(c^w)=m^w\pmod n.
\label{eq:rsa-homomorphism}
\end{equation}
رابطۀ دوم، ضرب $wm$ نیست. بنابراین نمی‌توان با توان‌رسانی متن رمز \en{RSA} همان مجموع وزن‌دار مورد نیاز را محاسبه کرد. همچنین \en{textbook RSA} قطعی و فاقد امنیت معنایی لازم برای رمزگذاری کاربردی است. طرح‌های امنی مانند \en{RSA-OAEP} لایۀ رمزگذاری متفاوتی دارند و خاصیت ضربی سادۀ نسخۀ پایه را برای این کاربرد حفظ نمی‌کنند.

\section{مقایسۀ نظری پنج طرح}
\begin{table}[htbp]
\centering
\caption{فضای پیام و نوع همریختی؛ تمام رابطه‌ها در پیمانۀ فضای پیام تفسیر می‌شوند}
\label{tab:he-theory}
\small
\begin{tabular}{llll}
\toprule
طرح & سال & فضای پیام & نتیجۀ ضرب متن‌های رمز \\
\midrule
\en{ElGamal} & 1985 & گروه $G$ & $m_1m_2$ \\
\en{Benaloh} & 1994 & $\Z_r$ & $m_1+m_2$ \\
\en{Paillier} & 1999 & $\Z_n$ & $m_1+m_2$ \\
\en{Damgård\lr{--}Jurik} & 2001 & $\Z_{n^s}$ & $m_1+m_2$ \\
\en{textbook RSA} & 1978 & $\Z_n$ & $m_1m_2$ \\
\bottomrule
\end{tabular}
\end{table}
انتخاب طرح به تابع مورد نیاز وابسته است. برای میانگین وزنیِ خروجی مدل، یک جمع وزن‌دار لازم است؛ ازاین‌رو پاییه، بنالوح و تخصص جمعی دامگارد\lr{--}یوریک تناسب جبری دارند. برای عملیات ضربی، الجمال و صورت پایۀ \en{RSA} رفتار متفاوتی ارائه می‌کنند. هیچ‌یک از این خواص به‌تنهایی گواه امنیت عملی کل سامانه نیست. سطح امنیت، اصالت پارامترها، مدیریت کلید، مدل مهاجم و صحت پیاده‌سازی نیز باید بررسی شوند.

\section{ارزیابی تجربی پیاده‌سازی‌های موجود}
\label{sec:crypto-benchmark}
منبع عددی این بخش فایل \en{Crypto/results/benchmark\_all.csv} است. این فایل ۱۵ ردیف، شامل پنج طرح در سه اندازۀ ۲۵۶، ۵۱۲ و ۱۰۲۴ بیت دارد. نتایج، زمان‌های همان اجرای ذخیره‌شده‌اند؛ اندازه‌گیری جدید روی سخت‌افزار یکسان یا میانگین چند تکرار برای این فایل ادعا نمی‌شود. واحد تمام زمان‌های این بخش میلی‌ثانیه است.

\subsection{صحت‌سنجی و تفاوت بار کاری}
در هر ۱۵ ردیف، وضعیت \en{homomorphic\_test} برابر \en{PASS} و خطا صفر است. برای طرح‌های جمعی، آزمون شامل رمزگذاری ۱۰۰ مقدار با مقیاس چهاررقم اعشار، تجمیع با وزن صحیح و مقایسۀ میانگین بازسازی‌شده با مقدار آشکار است. برای \en{RSA} و الجمال، آزمون، ضرب دو مقدار ۷ و ۱۱ و بازیابی ۷۷ است. در این دو طرح، \en{weighted\_average} برابر \en{NOT\_APPLICABLE} است؛ این وضعیت شکست آزمون نیست و همچنین موفقیت در محاسبۀ میانگین وزنی را نشان نمی‌دهد.

بار کاری دو خانواده برابر نیست. مقدار «۱۰۰ کلاینت» در فرادادۀ ردیف‌های ضربی نباید به‌عنوان رمزگذاری و رمزگشایی ۱۰۰ پیام در آن ردیف‌ها تفسیر شود. زمان‌های رمزگذاری و رمزگشایی آن‌ها در \en{CSV} جداگانه اندازه‌گیری نشده‌اند. علامت خط تیره در جدول‌ها همین نبود اندازه‌گیری را بیان می‌کند و مقدار صفر نیست. مقایسۀ مستقیم زمان کل این دو نوع آزمون برای رتبه‌بندی هم‌ارز امنیت یا توان عملیاتی معتبر نیست.

\begin{table}[htbp]
\centering
\small
\caption{زمان‌های طرح‌های جمعی برای ۱۰۰ پیام و تجمیع وزن‌دار}
\label{tab:crypto-additive}
\resizebox{\textwidth}{!}{%
\begin{tabular}{llrrrrr}
\toprule
طرح & بیت & تولید کلید & رمزگذاری & تجمیع/آزمون & رمزگشایی & کل \\
\midrule
\en{Paillier} & 256 & 4.6771 & 119.9491 & 1.6760 & 1.2705 & 127.5727 \\
\en{Damgård--Jurik ($s=1$)} & 256 & 11.0162 & 119.6526 & 1.1588 & 0.9399 & 132.7675 \\
\en{Benaloh} & 256 & 12.2062 & 7.9511 & 0.5458 & 2.1023 & 22.8054 \\
\en{Paillier} & 512 & 50.3984 & 730.4061 & 4.3244 & 6.7511 & 791.8800 \\
\en{Damgård--Jurik ($s=1$)} & 512 & 32.5281 & 696.0407 & 6.3247 & 7.8368 & 742.7303 \\
\en{Benaloh} & 512 & 67.0070 & 23.4188 & 1.4752 & 5.4431 & 97.3441 \\
\en{Paillier} & 1024 & 292.1688 & 3819.3751 & 11.7069 & 43.3571 & 4166.6079 \\
\en{Damgård--Jurik ($s=1$)} & 1024 & 276.9640 & 3176.5008 & 11.9147 & 31.8519 & 3497.2314 \\
\en{Benaloh} & 1024 & 310.4154 & 45.1094 & 3.3150 & 16.9272 & 375.7670 \\
\bottomrule
\end{tabular}}
\end{table}

\begin{table}[htbp]
\centering
\small
\caption{زمان‌های آزمون ضرب دو پیام؛ رمزگذاری و رمزگشایی جداگانه اندازه‌گیری نشده‌اند}
\label{tab:crypto-multiplicative}
\resizebox{\textwidth}{!}{%
\begin{tabular}{llrrrrr}
\toprule
طرح & بیت & تولید کلید & رمزگذاری & تجمیع/آزمون & رمزگشایی & کل \\
\midrule
\en{RSA} & 256 & 3.6671 & \en{--} & 0.5340 & \en{--} & 4.2011 \\
\en{ElGamal} & 256 & 1965.2741 & \en{--} & 2.4888 & \en{--} & 1967.7629 \\
\en{RSA} & 512 & 64.2466 & \en{--} & 2.3814 & \en{--} & 66.6280 \\
\en{ElGamal} & 512 & 15297.7712 & \en{--} & 11.3695 & \en{--} & 15309.1407 \\
\en{RSA} & 1024 & 412.0680 & \en{--} & 10.0805 & \en{--} & 422.1485 \\
\en{ElGamal} & 1024 & 224079.3934 & \en{--} & 35.5111 & \en{--} & 224114.9045 \\
\bottomrule
\end{tabular}}
\end{table}


\subsection{هزینۀ تولید کلید و رمزگذاری}
در اجرای ذخیره‌شده، تولید کلید الجمال بخش بزرگی از زمان آن طرح را تشکیل می‌دهد. این مشاهده به روش تولید پارامتر و آزمون مولد در همین کد مربوط است؛ به معنای غیرقابل‌استفاده بودن تمام پیاده‌سازی‌های الجمال نیست. در کاربردی که گروه استاندارد از پیش اعتبارسنجی شده و چند بار استفاده می‌شود، سهم این هزینه می‌تواند متفاوت باشد. برای سایر طرح‌ها نیز تعداد دفعات تولید کلید و امکان بازاستفادۀ امن از آن بر زمان مؤثر هر پرسش اثر دارد.

\begin{figure}[htbp]
\centering\includegraphics[width=0.95\textwidth]{original/crypto-key-generation.jpeg}
\caption{زمان تولید کلید در اجرای ذخیره‌شده؛ مقایسه به همین پیاده‌سازی‌ها محدود است}
\label{fig:crypto-keygen}
\end{figure}

رمزگذاری هر کلاینت در طرح‌های جمعی چند توان‌رسانی پیمانه‌ای دارد. افزایش طول پیمانه معمولاً هزینۀ محاسبات و اندازۀ متن رمز را افزایش می‌دهد. تفاوت زمانی بنالوح و پاییه در این آزمایش با تفاوت پیمانه‌ها، توان‌ها و فضای پیام مرتبط است، اما یک اجرای محدود علت‌سنجی قطعی همۀ این عوامل را ممکن نمی‌کند. در بنالوح، انتخاب $r$ بر رمزگشایی نیز اثر مستقیم دارد و باید همراه با ظرفیت مجموع وزن سنجیده شود.

\begin{figure}[htbp]
\centering\includegraphics[width=0.95\textwidth]{original/crypto-encryption.png}
\caption{زمان رمزگذاری در خانوادۀ جمعی؛ ستون‌های اندازه‌گیری‌نشدۀ خانوادۀ ضربی با صفر جایگزین نمی‌شوند}
\label{fig:crypto-encryption}
\end{figure}

\subsection{تجمیع، رمزگشایی و زمان کل}
تجمیع جمعی با ضرب توان‌دار متن‌های رمز انجام می‌شود. تعداد کلاینت‌ها، طول پیمانه و اندازۀ وزن‌ها بر آن اثر دارند. یک رمزگشایی نهایی می‌تواند از رمزگشایی همۀ پیام‌ها کم‌هزینه‌تر باشد، اما این کاهش تعداد عملیات به‌تنهایی محرمانگی در برابر دارندۀ کلید خصوصی را اثبات نمی‌کند. در بنالوح، حل لگاریتم گسستۀ محدود، بخش مشخصی از رمزگشایی است؛ در پاییه، توان‌رسانی و اعمال تابع $L$ انجام می‌شود.

زمان کل در دادۀ طرح‌های جمعی جمع زمان تولید کلید، رمزگذاری، تجمیع و رمزگشایی است. برای نمونه، در کلید ۲۵۶بیتی، این زمان برای پاییه ۱۲۷٫۵۷۲۷، برای تخصص دامگارد\lr{--}یوریک ۱۳۲٫۷۶۷۵ و برای بنالوح ۲۲٫۸۰۵۴ میلی‌ثانیه است. زمان کم‌تر بنالوح در این تنظیم نباید مستقل از فضای پیام محدود، روش رمزگشایی و سطح امنیت تفسیر شود. شکل‌های~\ref{fig:crypto-total} و~\ref{fig:crypto-performance} داده‌های ثبت‌شده را نمایش می‌دهند.

\begin{figure}[htbp]
\centering\includegraphics[width=0.95\textwidth]{original/crypto-total-time.png}
\caption{زمان کل آزمون‌های ذخیره‌شده؛ بار کاری خانواده‌های ضربی و جمعی متفاوت است}
\label{fig:crypto-total}
\end{figure}
\begin{figure}[htbp]
\centering\includegraphics[width=\textwidth]{original/crypto-performance.png}
\caption{مقایسۀ چندمعیارۀ زمان پیاده‌سازی‌های موجود در نسخۀ مبنا}
\label{fig:crypto-performance}
\end{figure}

\subsection{تأثیر اندازۀ کلید و حدود نتیجه‌گیری}
اندازه‌های ۲۵۶، ۵۱۲ و ۱۰۲۴ بیت در این بنچمارک، پارامترهای آموزشی و پژوهشی‌اند و تأیید امنیت تولیدی محسوب نمی‌شوند. یک طول بیت یکسان نیز لزوماً امنیت معادل در تمام خانواده‌ها ایجاد نمی‌کند. انتخاب پارامتر عملی باید بر اساس استاندارد جاری، طول عمر محرمانگی و تحلیل مستقل انجام شود.

برای سنجش دقیق‌تر، زمان تولید پارامتر باید از زمان هر پرسش جدا، بار کاری دو خانواده یکسان‌سازی، تکرارها و سخت‌افزار ثبت و بازۀ عدم‌قطعیت گزارش شود. روش‌های بهینۀ توان‌رسانی، موازی‌سازی و استفاده از کتابخانۀ رمزنگاری ممیزی‌شده از مسیرهای توسعه‌اند. تعمیم دامگارد\lr{--}یوریک برای $s>1$ نیز پیش از گزارش کارایی به پیاده‌سازی صحیح و آزمون‌های مستقل نیاز دارد. در این پایان‌نامه، تناسب جبری بنالوح با جمع وزن‌دار، مبنای بررسی پروتکل احرازاصالت‌شدۀ فصل بعد است؛ رتبۀ زمانی این اجرای محدود، ادعای برتری امنیتی عمومی ایجاد نمی‌کند.
\FloatBarrier
